\scanpage{78}
\mspage{13}
\uline{Question} Ces flèches sont elles des iso
\deleted{\unsure{--}} et si oui, comment expliciter un
iso en sens inverse ? \deleted{Si}

Si la détermination conjecturale des
syst. admissibles de disques sur une
surface holomorphe \inserted{compacte} ponctuée (de type
hyperbolique) était correcte, on trouverait aisément
\deleted{des \unsure{foncteurs}} \inserted{flèches} en sens
inverse; Au niveau plus naïf des
foncteurs (entre les catégories-sections
de ces multiplicités) on trouverait
ceci : à l'aide d'un \deleted{\unsure{iso}}
\[
u_i\in T_{x'_i}^{\times}\wedge_{\mathbb C^\times}T_{x''_i}^{\times}
\]
ou ce qui revient au $=$, un
\[
w_i(=u_i^{-1})\in T_{x'_i}^{\vee\times}\wedge T_{x''_i}^{\vee\times}
\]
\deleted{choisissant \unsure{iso}} \inserted{i.e. un} ``accouplement dualisant''
\[
\deleted{T_{x'_i}^{\vee}\xrightarrow{\sim} T_{x''_i}}\quad T_{x_i}\otimes T_{x''_i}\longrightarrow\mathbb C
\]
choisissant un \deleted{\unsure{métriq}} disque conforme
$D'_i$ de $T_{x'_i}$, \deleted{\unsure{on en déd}} d'où disque
conforme polaire $D''_i$ de $T_{x''_i}$, on
trouverait \deleted{des} pour ce syst. $(D'_i)$ $(D''_i)$
des réalisations de ces disques comme
% archive sheet number: 80
